\documentclass[10pt,a4paper]{article}
\usepackage[margin=23mm]{geometry}
\usepackage{booktabs,amsmath,graphicx,microtype,xcolor,float}
\usepackage{pgfplots}
\pgfplotsset{compat=1.18}
\usepackage[colorlinks=true,urlcolor=blue,citecolor=blue,linkcolor=blue]{hyperref}
\title{Completion-Time Feedback for Prefill Budgets:\\A Controlled Single-GPU Serving Study}
\author{Session D --- independent experimental report}
\date{4 October 2026}
\hypersetup{pdftitle={Completion-Time Feedback for Prefill Budgets},pdfauthor={Session D}}
\raggedbottom
\begin{document}
\maketitle
\begin{abstract}
We implement an aggregate prefill-token cap in native single-GPU serving, preserving request order, admission rules, model precision and kernels. Real mixed request traces show a repeatable fixed-budget throughput--generation-stall tradeoff. In the development trace, increasing the cap from 512 to 2,048 improves output throughput by 6.6\%, while increasing request maximum-gap P95. A four-level controller using only the previous completed mixed-step duration mostly selects one level at moderate load. We compare it with fixed budgets, a decode-count rule, its dominant fixed level and a matched-mean cap, including a normal-memory PRO 6000 pressure extension. The frozen moderate-load confirmation favors a fixed cap. At high load, feedback has a small exploratory joint-goodput advantage over the tested constants, with lower raw throughput and longer tail waiting. The contribution is a reproducible, bounded empirical result rather than a claim that adaptive chunking is new or broadly beneficial.

\end{abstract}

\section{Question and scope}
Long prompt processing shares execution time with ongoing token generation in
continuous batching. We test whether changing only the aggregate number of prompt
tokens allowed per iteration improves complete-service outcomes over a tuned,
deployable fixed cap. Equal token counts need not imply equal iteration costs,
but this is a hypothesis to test, not a premise establishing adaptive superiority.
This report distinguishes an observed throughput--stall tradeoff, an informative
state signal, and an end-to-end improvement: none implies the next automatically.

\section{Runtime intervention}
An initial launch failed before any request was executed because an environment
default overrode the requested in-process mode; it is retained as a zero-measurement
startup failure. Explicitly disabling the default multiprocessing environment
flag restores the requested in-process configuration. The runtime is vLLM 0.26.0, serving OLMoE-1B-7B-0924-Instruct in BF16 on one NVIDIA RTX PRO 6000
Blackwell Server Edition (97,887 MiB reported memory). The engine has a fixed
4,096 total-token limit, 192 sequence slots and 4,096 maximum context length. The normal-capacity
engine uses the runtime memory profiler with a 90\% device-memory target; no
manual small-KV limit is imposed. Resolved KV capacity is recorded. Prefix caching, speculative decoding and asynchronous
scheduling are disabled for every arm. The FCFS admission/allocation logic,
victim selection, model weights, precision and operators are common to all arms.
The first experiment is resident serving with no KV connector; zero preemptions
are checked, so recovery behavior is outside the evaluated regime.

A process-local scheduler patch introduces one counter $P$ per iteration. For
the running and waiting request loops, it clips only work beyond the remaining
prompt-token allowance. The counter is decremented only after successful native
allocation. Waiting requests use the native computed-prefix position; running
requests use their current position. The patch neither sorts requests nor changes
decode priority. At runtime we check prompt work $\leq P$, no preemption, and
exactly one scheduled token for every decoder already resident before the step.
The original installed scheduler is not modified.

The candidate action set is $P\in\{256,512,1024,2048\}$. Additional fixed
controls use 1,072 to match its development mean cap and 3,840 as a stronger
upper-budget endpoint; both receive dedicated full-trace warmups. Fixed arms use one cap for the
whole service trace. The count-only rule uses 2,048 below eight active decoders
and 512 otherwise. The feedback rule uses the same levels and adds the elapsed
time of the most recent completed mixed step. Starting at 1,024, it decreases
one level when the step exceeds a target $T=24$ ms, increases one level below $0.7T$,
and otherwise holds. With no active decoder it uses 2,048. The target is chosen after the first fixed comparison, whose mixed steps
span roughly 14--28 ms, and is frozen before the long-document holdout.
Only completed observations are used; no expert routing, extra forward pass, learned model,
full shape scan, or synchronization is added. The positive minimum cap and the
finite workload prevent indefinite withholding of prefill service in this domain.
This is not a worst-case wall-clock guarantee under arbitrary arrivals.

\section{Experimental method}
Each development trace has 36 requests: 12 short GSM8K prompts with 384 generated
tokens each and 24 long QMSum prompts with 160 tokens each. Eight short requests
arrive at zero, four at 3 seconds. Long requests arrive in three waves beginning
at 0.6, 2.5 and 5 seconds, with eight requests per wave and 40 ms within-wave
spacing. Arrivals are a common absolute clock, not completions from another
policy. The development prompt total is 92,625 tokens and every arm must generate
8,448 tokens and drain all requests. EOS is ignored for this controlled
scheduling experiment; this does not evaluate natural termination or task quality.
The PRO 6000 extension also includes 12-request low-load and 160-request high-load
traces. The high-load trace combines 32 short requests and 128 long prompts with
512 output tokens each (within the model context limit), totaling 471,895 prompt
tokens and 77,824 output tokens. Actual resident count and KV occupancy are
measured; the label high load alone does not prove a capacity bottleneck.
A separate long-document source trace replaces the 24 QMSum prompts; the short
background prompts are shared and therefore this is not a fully disjoint holdout.

All four cap levels run the full development trace as warmup before measurement,
using one common engine and compilation domain. Request state drains between
episodes and prefix caching is disabled. The first fixed comparison is ordered
512, 2,048, 2,048, 512. The 1,024 and 256 fixed controls and both dynamic rules are then compared
in forward and reverse orders. Fixed 2,048 is selected by development joint
goodput before holdout. The later 3,840 endpoint is an explicitly labeled
additional exploratory control. Later selection and confirmation steps are identified
separately in the results. A common host lock protects the complete campaign;
the controller verifies no other GPU process before initialization.

Every step records actual prompt/decode work, pre-step decoder count and summed
decoder context, policy selection overhead, scheduler time, and the elapsed
\texttt{engine.step()} envelope. This envelope includes host and device work
through returned outputs; it is not a CUDA-event kernel time. There is no added
device synchronization. Every request retains planned arrival, actual submission,
first scheduling, completion, output IDs, and every host-observed output time.

Throughput uses all generated tokens (8,448 in development) divided by the full trace duration,
including waiting and drain. TTFT and completion latency are relative to planned
arrival. For each request, maximum generation gap is the maximum difference
between consecutive observed output timestamps, excluding TTFT. Repeated
timestamps are reported rather than interpreted as observable intra-return token
timing. The primary joint SLO requires TTFT $\leq4$ s, maximum generation gap
$\leq100$ ms, and completion latency $\leq20$ s. Request goodput is the count
passing all three divided by full trace duration; output goodput counts only
outputs of passing requests. These thresholds are fixed before observing results
and describe an experimental contract, not a claimed production requirement.

\section{Results}
All 57 measured episodes completed, covering 4,636 request completions and 1,974,016 generated tokens. No measured request is dropped from the denominator. All token accounting checks pass.

\subsection{The fixed-budget tradeoff is real}

On the development trace, fixed 2,048 improves output throughput by 6.60\% relative to 512. Request maximum-gap P95 rises from 19.10 to 28.69 ms. Both satisfy the same primary joint SLO for every request. Thus a finer gap tradeoff exists even where the declared SLO does not reward smaller chunks.

\begin{table}[H]\centering\small
\begin{tabular}{lrrrrrr}
\toprule
Policy & $n$ & Output/s & TTFT P95 & Mean flow & Gap P95 & Goodput/s \\
 & & tokens & seconds & seconds & ms & requests \\ \midrule
Fixed 256 & 2 & 951.1 & 1.454 & 3.942 & 17.34 & 4.053 \\
Fixed 512 & 2 & 1037.0 & 0.651 & 3.492 & 19.10 & 4.419 \\
Fixed 1024 & 2 & 1085.7 & 0.266 & 3.244 & 21.45 & 4.626 \\
Fixed 2048 & 2 & 1105.4 & 0.104 & 3.134 & 28.69 & 4.711 \\
Decode count & 2 & 1036.8 & 0.650 & 3.457 & 22.80 & 4.418 \\
Feedback & 2 & 1087.5 & 0.254 & 3.221 & 26.83 & 4.634 \\
\bottomrule\end{tabular}
\caption{Development trace; fixed selection and candidate construction use these observations. Values are means of complete-run statistics; Gap P95 is across per-request maximum observed gaps, not pooled token gaps.}
\end{table}

\subsection{Frozen long-document confirmation}

On the long-document holdout, feedback changes throughput by -0.89\% and joint goodput by -0.89\% versus the development-selected fixed 2,048. The short background prompts are shared with development. The three main-policy repetitions provide repeatability evidence, not a broad workload-population confidence interval.

\begin{table}[H]\centering\small
\begin{tabular}{lrrrrrr}
\toprule
Policy & $n$ & Output/s & TTFT P95 & Mean flow & Gap P95 & Goodput/s \\
 & & tokens & seconds & seconds & ms & requests \\ \midrule
Fixed 2048 & 3 & 1103.6 & 0.101 & 3.146 & 28.44 & 4.703 \\
Fixed 1024 & 2 & 1092.7 & 0.251 & 3.212 & 21.56 & 4.656 \\
Decode count & 3 & 1042.2 & 0.616 & 3.436 & 22.69 & 4.441 \\
Feedback & 3 & 1093.8 & 0.244 & 3.193 & 25.99 & 4.661 \\
\bottomrule\end{tabular}
\caption{Frozen long-document holdout. Values are means of complete-run statistics; Gap P95 is across per-request maximum observed gaps, not pooled token gaps.}
\end{table}

\subsection{Normal-memory pressure extension}

\begin{table}[H]\centering\small
\begin{tabular}{lrrrrrr}
\toprule
Policy & $n$ & Output/s & TTFT P95 & Mean flow & Gap P95 & Goodput/s \\
 & & tokens & seconds & seconds & ms & requests \\ \midrule
Fixed 1024 & 2 & 2114.7 & 11.349 & 25.913 & 66.74 & 0.652 \\
Fixed 2048 & 2 & 2232.1 & 4.760 & 27.263 & 70.54 & 0.688 \\
Decode count & 2 & 1867.6 & 22.391 & 25.228 & 44.37 & 1.092 \\
Feedback & 2 & 1600.9 & 32.452 & 25.438 & 32.07 & 1.121 \\
\bottomrule\end{tabular}
\caption{High load: 160 requests, normal memory allocation. Values are means of complete-run statistics; Gap P95 is across per-request maximum observed gaps, not pooled token gaps.}
\end{table}

\begin{table}[H]\centering\small
\begin{tabular}{lrrrrrr}
\toprule
Policy & $n$ & Output/s & TTFT P95 & Mean flow & Gap P95 & Goodput/s \\
 & & tokens & seconds & seconds & ms & requests \\ \midrule
Fixed 2048 & 4 & 476.1 & 0.118 & 1.786 & 25.15 & 2.029 \\
Feedback & 2 & 476.5 & 0.158 & 1.806 & 24.87 & 2.030 \\
\bottomrule\end{tabular}
\caption{Low load: 12 requests. Decode-count and fixed 2,048 choose exactly the same cap at every step and are merged here. Values are means of complete-run statistics; Gap P95 is across per-request maximum observed gaps, not pooled token gaps.}
\end{table}

\subsection{Dominant-level and mean-budget controls}

In development, feedback selects 1,024 in 84 of 88 mixed steps in each repetition, with four selections of 2,048. Its mean mixed cap is 1,070.55. Fixed 1,024 is the direct dominant-level control; fixed 1,072 additionally approximates the mean allowance. This comparison counts full realized service time, rather than assuming identical state trajectories from similar average caps.

\begin{table}[H]\centering\small
\begin{tabular}{lrrrrrr}
\toprule
Policy & $n$ & Output/s & TTFT P95 & Mean flow & Gap P95 & Goodput/s \\
 & & tokens & seconds & seconds & ms & requests \\ \midrule
Fixed 1072 & 2 & 1093.5 & 0.239 & 3.214 & 21.81 & 4.660 \\
Fixed 3840 & 2 & 1103.3 & 0.057 & 3.152 & 39.69 & 4.701 \\
Feedback & 2 & 1087.0 & 0.256 & 3.227 & 26.95 & 4.632 \\
\bottomrule\end{tabular}
\caption{Additional exploratory controls: matched mean cap and upper fixed endpoint. Values are means of complete-run statistics; Gap P95 is across per-request maximum observed gaps, not pooled token gaps.}
\end{table}

\begin{table}[H]\centering\small
\begin{tabular}{lrrrrrr}
\toprule
Policy & $n$ & Output/s & TTFT P95 & Mean flow & Gap P95 & Goodput/s \\
 & & tokens & seconds & seconds & ms & requests \\ \midrule
Fixed 256 & 2 & 1513.5 & 35.271 & 27.227 & 28.36 & 0.778 \\
Fixed 304 & 2 & 1622.1 & 31.387 & 26.117 & 38.50 & 0.896 \\
Fixed 384 & 2 & 1748.5 & 26.934 & 25.315 & 35.46 & 1.078 \\
Fixed 512 & 2 & 1858.2 & 22.577 & 25.437 & 77.30 & 0.825 \\
Fixed 768 & 2 & 2028.5 & 15.724 & 25.455 & 61.65 & 0.834 \\
Fixed 3840 & 2 & 2273.5 & 2.058 & 28.029 & 83.14 & 0.000 \\
Feedback & 2 & 1605.9 & 32.336 & 25.381 & 32.68 & 1.125 \\
\bottomrule\end{tabular}
\caption{Additional high-load fixed controls; constants apply to the entire trace, never individual segments. Values are means of complete-run statistics; Gap P95 is across per-request maximum observed gaps, not pooled token gaps.}
\end{table}

Across measured episodes, peak resident concurrency is 160 and peak observed KV occupancy is 88.31\%. All measured preemption counts are zero. The largest per-run P95 policy-choice time is 3.56 microseconds; this excludes common observation construction, which remains included in all end-to-end times. No second model is introduced because an independently useful extra signal has not yet been established.

\subsection{A conditional gain, not a uniform improvement}

In the additional high-load comparison, feedback reaches 1.125 qualified requests/s versus 1.078 for fixed 384, the highest mean among the tested high-load fixed constants. This is 4.28\% higher goodput, alongside 8.15\% lower raw throughput, 20.06\% higher TTFT P95, and 9.52\% higher completion-latency P95. This is a measured tradeoff, not Pareto dominance. Fixed 304 approximates feedback's 303.42-token realized mixed-step mean and 304.03-token mean cap in this supplementary group, but has only 0.896 qualified requests/s. Mean work alone does not explain the high-load goodput difference.

The original high-load main group has zero gap-SLO failures. Feedback admits 30--31 first-wave long requests into the joint SLO versus 21--22 for the count rule; later waves fail under both. Compared with count-only scheduling, feedback improves complete-request mean outcomes for 87 of 160 request IDs and worsens 73; the joint-goodput increment comes only from early requests crossing thresholds. Third-wave mean completion rises from 34.57 to 39.62 seconds. This explicitly exposes the service cost of the gain.

All fixed-512 supplementary trials remain in the report: its goodput ranges from 0.571 to 1.078 requests/s. In the slower-SLO trial, an observed 108.97 ms gap causes 67 gap failures; the associated step envelope is 108.47 ms and scheduler time is 1.07 ms. This record is neither removed nor attributed to a specific GPU or host cause. Fixed-3840 goodput is zero because all 160 requests exceed the 20-second completion constraint, despite the largest raw throughput; none fails its gap condition. These results show sensitivity to the full SLO conjunction, not only average step latency.

The complete retained log contains one monitored Triton JIT warning during the initial 512 warmup and no recorded JIT compilation warning in the measured episodes. This is a log-based check, not a claim that every source of runtime jitter is observable.


\subsection{What the observed state explains}
At exactly 512 actual prompt tokens, mean completed-step time rises from
13.74 to 17.81 ms when the decoder-count bin changes from 1--7 to 16--23;
the reverse repetition gives 14.07 to 17.93 ms. At 2,048 actual prompt tokens,
the 8--15 versus 16--23 bins give 24.92 versus 27.38 ms, reproduced as
24.93 versus 27.37 ms. These comparisons hold prompt work fixed, not total tokens.
For a stricter fixed-total comparison, two fixed-1,024 runs each contain 19 steps
with exactly 10 decoders. Splitting them by context sum gives mean costs
18.268 versus 18.730 ms and 18.281 versus 18.741 ms. This is a small repeatable
association at 1,034 total scheduled tokens. A five-step bin with 18 decoders
has the opposite direction, so we do not infer a monotonic causal law or a
validated independent prediction signal.

\begin{figure}[htbp]\centering
\begin{tikzpicture}
\begin{axis}[width=.82\textwidth,height=5.3cm,xlabel={Active decoders},
 ylabel={Completed step envelope (ms)},grid=major,
 legend style={at={(.03,.97)},anchor=north west},legend cell align=left]
\addplot[only marks,mark=*,mark size=1.4pt,blue] table[x=decode,y=ms,col sep=comma]{step_cost_512.csv};
\addlegendentry{Exactly 512 prefill tokens}
\addplot[only marks,mark=triangle*,mark size=1.8pt,red] table[x=decode,y=ms,col sep=comma]{step_cost_2048.csv};
\addlegendentry{Exactly 2,048 prefill tokens}
\end{axis}
\end{tikzpicture}
\caption{Actual mixed steps from the first measured fixed pair. Partial tail
chunks are excluded only from this equal-prompt-work diagnostic, not from service metrics.}
\end{figure}

The fixed-1,024 and feedback development runs agree on the complete output
sequence for 24 of 36 requests; each policy agrees on all 36 sequences between
its own repetitions. Every run still generates exactly 8,448 tokens. This is
consistent with shape-dependent numerical execution changing greedy trajectories;
it does not establish either a precision change or a task-quality difference.


\section{Related work and interpretation}
Sarathi-Serve establishes chunked prefills and decode-maximal batches as tools for
the throughput--latency tradeoff~\cite{sarathi}. Its current official implementation
also supports decreasing chunk sizes with progress through a prompt~\cite{sarathicode}.
QoServe chooses chunks using deadline state and predicted execution cost from
request count and context information~\cite{qoserve}. These are direct prior art,
so neither adaptive chunking nor context awareness is claimed here as new.
DLFP is particularly close: it uses proportional latency feedback to bound
overlapping prefill, and reports both small-model gains and failed transfer to
larger configurations~\cite{dlfp}. Its asynchronous scheduler-cadence measurement
and unsuccessful completion-aware revision make measurement fidelity and strong
fixed controls central to this experiment. A synchronous host completion
envelope is a practical observation choice here, not by itself a novelty claim.

\section{Limitations and conclusion}
These experiments establish the fixed-budget tradeoff and a conditional high-load goodput gain, but not a consistent improvement across full service metrics and loads. The frozen long-document confirmation does not beat its development-selected fixed policy. At moderate load, most actions collapse to the same cap; at higher load, the feedback target can be incompatible with the decode cost already present, so reducing prompt service can move delay into the waiting queue. The scope is one model, one GPU, in-process synchronous vLLM, fixed output lengths, and a small set of structured arrival traces. Model precision is unchanged, but different batch shapes can alter greedy outputs through floating-point execution; fixed output lengths hold service work constant and are not a semantic-equivalence claim. Host-observed delivery gaps omit network transport and do not isolate device-only latency. GPU clocks are not locked. A 100 ms gap SLO leaves many small-load requests qualified under every policy, so its goodput is often determined by full service duration. No hindsight choice of a fixed cap per trace segment is presented as an online algorithm. The high-load trace was used for fixed-budget tuning and is exploratory rather than an independent high-pressure holdout. The count-only comparator is a conservative simple rule, not a tuned count-based controller. Thus the value of elapsed time beyond a strong count rule remains unresolved; a MoE label supplies no novelty claim.


\clearpage
\section*{Artifact and reproduction}
The accompanying artifact contains the process-local scheduler intervention,
the serving driver, frozen tokenized workloads, exact executed phase commands,
per-step and per-request originals, numerical analysis, and this paper's sources.
The \texttt{reproduction/} directory is a clean execution bundle: its command
files are reconstructed from the retained executed protocols, including the
warmup-only phases. The repository README gives the launch and analysis commands.
Replaying the campaign requires the recorded vLLM/Torch environment and the
model weights; the artifact does not package weights or credentials.

The resolved stack is Torch 2.11.0+cu130, CUDA 13.0, vLLM 0.26.0 and NVIDIA
driver 595.71.05. The model revision is
\texttt{7f1c97f440f06ce36705e4f2b843edb5925f4498}; all three weight-shard SHA-256
digests match the retained model manifest. Native and patched scheduler digests
are recorded in the initial phase's \texttt{patch.json}.
The memory profiler allocated 71.81 GiB of KV cache, with 36,765 blocks of
16 tokens (one block is reserved). This is normal device-capacity allocation,
not a reduced-cache mechanism experiment.

The analyzer validates request count, completion, fixed output length, prompt
and decode token conservation, timestamp counts, and prefill budget compliance.
The report builder excludes warmups and the failed zero-request startup from
its 57 measured episodes. All 108 archived remote originals were verified
byte-for-byte after transfer. The campaign released its shared GPU lock after
the last comparison; analysis and document preparation continued locally.

\begin{thebibliography}{9}
\bibitem{sarathi} A. Agrawal et al. Taming Throughput-Latency Tradeoff in LLM
Inference with Sarathi-Serve. OSDI, 2024.
\url{https://arxiv.org/abs/2403.02310}
\bibitem{sarathicode} Microsoft. Sarathi-Serve scheduler source, accessed 4 October 2026.
\url{https://github.com/microsoft/sarathi-serve/blob/main/sarathi/core/scheduler/sarathi_scheduler.py}
\bibitem{qoserve} QoServe: Breaking the Silos of LLM Inference Serving. ASPLOS, 2026.
\url{https://www.microsoft.com/en-us/research/wp-content/uploads/2025/10/QoServe_CC_final_v2.pdf}
\bibitem{dlfp} G. Agarwal, A. Garg and I. Singhal. Decode-Latency Feedback Prefill:
A Model-Free Controller and Its Generalization Limits. 2026.
\url{https://arxiv.org/abs/2609.38386}
\end{thebibliography}
\end{document}
